\SourceLeaf{300}
\begin{flushright}
Göttingen, 23. Juli 1808.
\end{flushright}

\par\indent Auf Ew. --- Wunsch teile ich Ihnen hier die Beobachtungen mit, die ich seither an den neuen Planeten machen konnte. Sie wurden alle mit dem Kreismikrometer gemacht. Mit unserem Mauerquadranten wird dieses Jahr nur Vesta beobachtet werden können, und sie kulminiert noch bei Tage.

\Needspace{9\baselineskip}
\begin{center}\emph{Beobachtungen der Juno.}\end{center}

\begin{center}\fontsize{10}{12}\selectfont\setlength{\tabcolsep}{3pt}\renewcommand{\arraystretch}{1.15}\begin{tabular}{lrrr}
1808 & \shortstack{Mittl. Zeit\\in Göttingen} & \shortstack{Rektaszension} & \shortstack{südliche\\Deklination}\\
Juni 20 & $11^{\text{u}}\,49^{\text{m}}\,0^{\text{s}}$ & $315^\circ\;29'\;34''$ & $2^\circ\;16'\;23''$\\
22 & $12\;0\;45$ & $315\;23\;1$ & $2\;14\;29$\\
Juli 6 & $12\;42\;23$ & $314\;1\;0$ & $2\;19\;7$\\
\end{tabular}\end{center}

\Needspace{9\baselineskip}\begin{center}\emph{Beobachtungen der Vesta.}\end{center}

\begin{center}\fontsize{9}{12}\selectfont\setlength{\tabcolsep}{3pt}\renewcommand{\arraystretch}{1.15}\begin{tabular}{lrrrl}
Juni 22 & $13^{\text{u}}\,46^{\text{m}}\,25^{\text{s}}$ & $353^\circ\;12'\;17''$ & $9^\circ\;12'\;16''$ & \\
Juli 1 & $14\;38\;57$ & $354\;28\;3$ & $9\;9\;39$ & \\
2 & $12\;50\;2$ & $355\;7\;7$ & & (von Prof. \textsc{Harding}).\\
13 & $13\;22\;37$ & $356\;35\;27$ & $9\;27\;40$ & \\
\end{tabular}\end{center}

\Needspace{7\baselineskip}\begin{center}\emph{Beobachtungen der Pallas.}\end{center}

\begin{center}\fontsize{9}{12}\selectfont\setlength{\tabcolsep}{3pt}\renewcommand{\arraystretch}{1.15}\begin{tabular}{lrrrl}
Juli 13 & $12^{\text{u}}\,3^{\text{m}}\,17^{\text{s}}$ & $300^\circ\;21'\;49''$ & $18^\circ\;41'\;44''$ nördl. & \\
15 & $12\;1\;28$ & $299\;57\;3$ & $18\;32\;40$ & (von Prof. \textsc{Harding}).\\
\end{tabular}\end{center}

\par\indent Vesta hat jetzt schon die $7^{\text{te}}$ Größe. Juno und Pallas sind dagegen sehr lichtschwach und schätzungsweise kaum von $10^{\text{ter}}$ Größe. Ceres habe ich noch nicht gesehen. Die Beobachtungen der Pallas werden wenig zuverlässig sein, da sie nur mit einem schwächeren Instrument gemacht werden konnten, weil der Planet für das Herschel-Teleskop jetzt zu hoch steht. --- Wegen der Elemente verweise ich auf Nr. 14, 40 und 107 unserer gelehrten Anzeigen [Seite 292, 299, 300 dieses Bandes]. Ephemeriden für das nächste Jahr werde ich erst nach einer erneuten Verbesserung der Elemente erstellen. Bei Vesta halte ich eine Ephemeride umso weniger für nötig, da sie, wenn sie im März 1809 unsichtbar geworden ist, erst im Oktober oder November 1809 wieder sichtbar wird. --- Hier noch einige Eintritte der Jupitermonde:

\begin{center}
\fontsize{10}{12}\selectfont\setlength{\tabcolsep}{3pt}\begin{tabular}{lccl}
$2^{\text{ter}}$ Mond Juni 22. & $12^{\text{u}}\,45^{\text{m}}\,9^{\text{s}}$ & & \textsc{Harding} $3\tfrac12$f. \textsc{Dollond}.\\
& $12\;45\;32$ & & \textsc{Gauss} $10$f. \textsc{Herschel}.\\
$1^{\text{ster}}$ Mond Juli 13. & $13\;45\;36$ & & \textsc{Harding} $4$f. \textsc{Dollond}.\\
& $13\;45\;38$ & & \textsc{Thiarks} $3\tfrac12$f. \textsc{Dollond}.\\
& $13\;45\;38$ & & \textsc{Gauss} $10$f. \textsc{Herschel}.\\
\end{tabular}

Austritte.

\fontsize{10}{12}\selectfont\setlength{\tabcolsep}{3pt}\begin{tabular}{lccl}
$4^{\text{ter}}$ Mond Juli 14. & $13\;35\;16$ & & \textsc{Harding} $10$f. \textsc{Herschel}.\\
& $13\;35\;56$ & & \textsc{Thiarks} $3\tfrac12$f. \textsc{Dollond}.\\
$2^{\text{ter}}$ Mond Juli 24. & $12\;32\;13$ & & \textsc{Thiarks} $3\tfrac14$f. \textsc{Dollond}.\\
& $12\;32\;27$ & & \textsc{Harding} $4$f. \textsc{Dollond}.\\
\end{tabular}
\end{center}

\Needspace{8\baselineskip}\par\indent Außerdem 2 Sternbedeckungen.

\smallskip
\noindent Juni 4.\quad $9^{\text{u}}\,3^{\text{m}}\,39^{\text{s}}0$\quad Eintritt von $\iota$ Virginis. \textsc{Gauss} und \textsc{Harding} in derselben Sekunde.

\smallskip
\noindent Juli 6.\quad $10\;39\;6.2$\quad Eintritt von $\mu$ Sagittarii. \textsc{Gauss}.

\smallskip
{\rightskip=0pt plus 3em\par\indent Der König hat den Instrumentenbestand unserer Sternwarte vor Kurzem durch ein Geschenk erweitert, ein Spiegelteleskop von \textsc{Chevalier} und einen Zeitmesser von \textsc{Berthoud}. --- Der Druck meines
\SourceLeaf{301}Werks über die Bestimmung der Planetenbahnen ist noch nicht abgeschlossen. Bisher sind erst 12 Bogen fertig, und das Ganze wird etwa 30 umfassen. Wir erwarten das $4^{\text{te}}$ Blatt der Karten von Professor \textsc{Harding} mit jeder Post. Wenn dieses wertvolle Unternehmen doch beim Publikum genügend Unterstützung fände!~---\par}
